\section{Server Recovery Strategies}
\label{sec:server-recovery}

Workflow schedulers (cron and poll triggers) live inside BullMQ, which persists its state in Redis. Postgres, however, remains the single 
source of truth for \emph{which} projects are supposed to be deployed. Whenever the Node.js process shuts down — deliberately or by 
crashing — these two stores can drift apart: Redis may still hold scheduler entries for workflows that Postgres no longer considers active, 
or Postgres may say a project is \texttt{deployed: true} while no scheduler for it actually exists anymore. The recovery subsystem exists 
to reconcile this drift, and it does so differently depending on whether the shutdown was intentional (development) or unexpected 
(production).

Two functions implement this behaviour in \texttt{workflow-runner.ts}: \texttt{gracefulStop()}, which runs on a clean local shutdown, 
and \texttt{recoverDeployedWorkflows()}, which runs on server boot in production.

\subsection{Two Recovery Paths}

The system deliberately treats development and production differently, because the two environments fail in different ways.

\begin{table}[H]
\centering
\caption{Recovery strategy by environment}
\label{tab:recovery-strategy}
\begin{tabular}{|p{2.8cm}|p{5.2cm}|p{5cm}|}
\hline
\textbf{Environment} & \textbf{Trigger} & \textbf{Strategy} \\
\hline
Development & \texttt{SIGINT}/\texttt{SIGTERM}/\texttt{SIGTSTP} on local shutdown & \texttt{gracefulStop()} — actively tears down every 
scheduler before the process exits \\
\hline
Production & Server startup (\texttt{index.ts}) & \texttt{recoverDeployedWorkflows()} — assumes the previous process died without cleanup 
and resyncs from Postgres \\
\hline
\end{tabular}
\end{table}

Both functions guard themselves against running in the wrong environment:

\begin{lstlisting}
export const gracefulStop = async () => {
  if (process.env.ENV === 'production') return;
  ...
};

export const recoverDeployedWorkflows = async () => {
  if (process.env.ENV === 'development') return;
  ...
};
\end{lstlisting}

In development, restarting the server (e.g. via \texttt{nodemon}) is frequent and expected, so schedulers are stopped cleanly on every 
shutdown — there is no need to \emph{recover} anything on the next boot, because nothing was left dangling. In production, a crash is the 
failure mode being planned for: the process could die at any point without running its shutdown hooks, so recovery has to assume Redis and 
Postgres have already diverged and must be resynced from scratch when the server comes back up.

\subsection{Dev Mode: \texttt{gracefulStop()}}

On a local shutdown, the function queries Postgres for every project that is currently marked \texttt{deployed: true} and 
\texttt{active: true}, rebuilds each one's workflow graph, and calls \texttt{stopTriggerNodes()} on it to remove the corresponding BullMQ 
schedulers:

\begin{lstlisting}
const deployedProjects = await prisma.project.findMany({
  where: { deployed: true, active: true },
  select: { id: true, workflow: true }
});

for (const project of deployedProjects) {
  await stopTriggerNodes(buildWorkflowGraph(project), project.id);
}
\end{lstlisting}

Once every scheduler has been removed, it performs a single atomic update: any execution still marked \texttt{NEW} or \texttt{RUNNING} is 
transitioned to \texttt{CANCELLED} with an explanatory error message, and every affected project is flipped back to \texttt{deployed: false}:

\begin{lstlisting}
await prisma.$transaction([
  prisma.execution.updateMany({
    where: {
      projectId: { in: projectIds },
      status: { in: [ExecutionStatus.NEW, ExecutionStatus.RUNNING] }
    },
    data: {
      status: ExecutionStatus.CANCELLED,
      endedAt: new Date(),
      error: 'Server shutdown - execution canceled'
    }
  }),
  prisma.project.updateMany({
    where: { id: { in: projectIds } },
    data: { deployed: false }
  })
]);
\end{lstlisting}

Wrapping both updates in a single \texttt{\$transaction} matters here: it guarantees that Postgres never ends up in a state where an execution 
is left \texttt{RUNNING} while its parent project is already marked \texttt{deployed: false}, or vice versa. Since this path runs before the process actually exits, the developer gets a workspace that is fully clean the next time the server starts — no stale schedulers in Redis, no dangling executions in Postgres.

\subsection{Production: \texttt{recoverDeployedWorkflows()}}

Production cannot rely on a clean shutdown sequence running, so recovery happens the other way around: on startup, before the server accepts
 new deployments, this function walks through what the \emph{previous} process left behind and repairs it in three phases.

\subsubsection{Phase 1 --- Clean up dangling executions}

Any execution still sitting in \texttt{NEW} or \texttt{RUNNING} status could only have gotten there if the process died mid-flight; a healthy
 shutdown would have already cancelled them (dev) or nothing would have started a new one this early (prod). These are reclassified as \texttt{CRASHED} 
 rather than \texttt{CANCELLED}, since the distinction is diagnostically useful — a cancellation is intentional, a crash is not:

\begin{lstlisting}
const cleanupResult = await prisma.execution.updateMany({
  where: { status: { in: [ExecutionStatus.NEW, ExecutionStatus.RUNNING] } },
  data: {
    status: ExecutionStatus.CRASHED,
    endedAt: new Date(),
    error: 'Server crashed unexpectedly - execution terminated'
  }
});
\end{lstlisting}

\subsubsection{Phase 2 --- Reload deployed projects}

Postgres is treated as ground truth here: any project still flagged \texttt{deployed: true} and \texttt{active: true} is assumed to need its 
triggers running again, regardless of what state Redis happens to be in.

\begin{lstlisting}
const deployedProjects = await prisma.project.findMany({
  where: { deployed: true, active: true },
  select: { id: true, workflow: true }
});
\end{lstlisting}

\subsubsection{Phase 3 --- Resync, don't restore}

For each deployed project, the function does not assume any existing scheduler is trustworthy. It first calls \texttt{stopTriggerNodes()} to 
remove whatever might already be registered under that project's deterministic job IDs (\texttt{cron-\{projectId\}-\{nodeId\}}, \texttt{poll-\{projectId\}-\{nodeId\}}), 
and only then calls \texttt{enqueueTriggerNodes()} to schedule them fresh:

\begin{lstlisting}
for (const project of deployedProjects) {
  const workflowGraph = buildWorkflowGraph(project);

  try {
    await stopTriggerNodes(workflowGraph, project.id);
  } catch (err) {
    console.error(`Failed to remove stale schedulers for ${project.id}`, err);
  }

  await enqueueTriggerNodes({
    projectId: project.id,
    workflowGraph: workflowGraph
  });
}
\end{lstlisting}

This stop-then-reenqueue sequence is a deliberate design choice, not an oversight: it is a \emph{resync}, not a \emph{restore}. If \texttt{upsertJobScheduler} 
were relied upon alone, a workflow definition that changed while the server was down (e.g. a node's interval parameter was edited) could leave BullMQ 
holding a scheduler with a mixture of old and new configuration. Explicitly removing first guarantees that whatever gets registered reflects
 exactly what is currently defined in Postgres, with no possibility of drift surviving the recovery pass.

Each project is recovered independently inside its own \texttt{try/catch}, so a single malformed workflow graph cannot block the rest of the fleet from redeploying:

\begin{lstlisting}
try {
  ...
} catch (err) {
  console.error(`[Recovery Error] Failed to recover project ${project.id}:`, err);
}
\end{lstlisting}

\subsection{Where Recovery Is Wired In}

\texttt{recoverDeployedWorkflows()} is invoked once, from the HTTP server's startup callback in \texttt{index.ts}, strictly gated behind the production check:

\begin{lstlisting}
const server = app.listen(PORT, () => {
  console.log(`Server is running on port ${PORT}`);

  if (process.env.ENV === 'production') {
    recoverDeployedWorkflows()
      .then(() => console.log('Recovery initialization triggered successfully.'))
      .catch((error) =>
        console.error('Critical failure triggering workflow recovery on startup:', error)
      );
  }
});
\end{lstlisting}

\texttt{gracefulStop()} is invoked from the same file's signal handlers, but only on the development branch of an \texttt{if/else} that also handles 
the production shutdown path (which cancels executions directly via Prisma but deliberately leaves \texttt{deployed: true} so that the next production
 boot's recovery pass has something to act on):

\begin{lstlisting}
if (process.env.ENV === 'production') {
  await prisma.$transaction([
    prisma.execution.updateMany({
      where: { status: { in: [ExecutionStatus.NEW, ExecutionStatus.RUNNING] } },
      data: {
        status: ExecutionStatus.CANCELLED,
        endedAt: new Date(),
        error: 'Server shutdown - execution canceled'
      }
    })
    recoverDeployedWorkflows() // relies on it after restart
  ]);
} else {
  await gracefulStop();
}
\end{lstlisting}

This is the key asymmetry between the two paths: in development, \texttt{deployed} is reset to \texttt{false} because the developer expects 
the workflow to genuinely stop. In production, \texttt{deployed} is left as-is on shutdown, because a production shutdown is far more likely 
to be a deploy/restart than a deliberate stop, and the intent is for the workflow to come back up automatically.

\subsection{Postgres as the Source of Truth}

Both recovery paths converge on the same underlying principle: Redis and BullMQ hold \emph{ephemeral} scheduling state, while Postgres holds 
the \emph{durable} record of what a workflow's deployment state ought to be. Neither function ever infers a project's intended deployment status 
by inspecting BullMQ — it always starts from a Postgres query. This means that even in the worst case, where Redis has been flushed or a job 
scheduler has silently disappeared, recovery on the next production boot will still reconstruct the correct set of running workflows purely 
from the \texttt{deployed} flag on each \texttt{Project} row.